\documentclass[12pt,a4paper]{article}
\usepackage[a4paper,top=12mm,bottom=14mm,left=15mm,right=15mm]{geometry}
\usepackage[T1]{fontenc}
\usepackage{amsmath}
\usepackage{amssymb}
\usepackage{enumitem}
\usepackage{multicol}
\usepackage{array}
\usepackage{xcolor}
\usepackage[most]{tcolorbox}

\pagestyle{plain}
\setlength{\parindent}{0pt}
\setlength{\parskip}{4pt}
\setlist[enumerate,1]{label=\textbf{(\alph*)},leftmargin=*,itemsep=3pt,topsep=3pt,parsep=1pt}
\setlist[enumerate,2]{label=\roman*),leftmargin=*,itemsep=2pt,topsep=2pt,parsep=0pt}

\newtcolorbox{solution}[1]{breakable,colback=blue!1,colframe=blue!45!black,boxrule=0.6pt,
  colbacktitle=blue!8,coltitle=blue!50!black,
  left=6pt,right=6pt,top=4pt,bottom=4pt,
  title={\bfseries #1},fonttitle=\normalsize,
  toptitle=1pt,bottomtitle=1pt}

\begin{document}

\begin{center}
  {\Large\bfseries The Laws of Indices}\\[2pt]
  {\large Worked Solutions}
\end{center}

\begin{tcolorbox}[colback=orange!2,colframe=orange!60!black,boxrule=0.6pt]
\textbf{Remember:} For powers of the same base,
\[
a^m a^n=a^{m+n},\qquad \frac{a^m}{a^n}=a^{m-n},\qquad (a^m)^n=a^{mn},
\]
and \(a^0=1\), \(a^{-n}=1/a^n\). A negative index takes a reciprocal; it does not automatically make the answer negative.
\end{tcolorbox}

\begin{solution}{Q1. Multiplying powers}
Use \(a^m\times a^n=a^{m+n}\).
\begin{enumerate}[label=\textbf{(\alph*)}]
\item \(5^2\times5^4=5^{2+4}=5^6\).
\item \(x^3\times x^6=x^{3+6}=x^9\).
\item \(3^5\times3^2=3^{5+2}=3^7\).
\item \(10^4\times10^3=10^{4+3}=10^7\).
\item \(7\times7^4=7^1\times7^4=7^{1+4}=7^5\).
\item \(y^2\times y^9=y^{2+9}=y^{11}\).
\item \(m\times m^5=m^1\times m^5=m^{1+5}=m^6\).
\item \(a^2\times a\times a^6=a^{2+1+6}=a^9\).
\item \(6^7\times6^3=6^{10}\), since \(7+3=10\). The missing index is \(7\).
\item \(b^4\times b^7=b^{11}\), since \(4+7=11\). The missing index is \(7\).
\end{enumerate}
\end{solution}

\begin{solution}{Q2. Dividing powers}
Use \(a^m\div a^n=a^{m-n}\).
\begin{enumerate}[label=\textbf{(\alph*)}]
\item \(6^7\div6^3=6^{7-3}=6^4\).
\item \(\dfrac{x^9}{x^4}=x^{9-4}=x^5\).
\item \(5^8\div5^2=5^{8-2}=5^6\).
\item \(y^{10}\div y^3=y^{10-3}=y^7\).
\item \(\dfrac{9^6}{9}=\dfrac{9^6}{9^1}=9^{6-1}=9^5\).
\item \(\dfrac{a^7}{a^6}=a^{7-6}=a^1=a\).
\item \(\dfrac{3^4\times3^5}{3^2}=\dfrac{3^{4+5}}{3^2}=\dfrac{3^9}{3^2}=3^{9-2}=3^7\).
\item \(\dfrac{p^3\times p^8}{p^2\times p^4}=\dfrac{p^{3+8}}{p^{2+4}}=\dfrac{p^{11}}{p^6}=p^{11-6}=p^5\).
\item \(10^9\div10^4=10^{9-4}=10^5\). The missing index is \(9\).
\item \(\dfrac{b^{12}}{b^9}=b^{12-9}=b^3\). The missing index is \(9\).
\end{enumerate}
\end{solution}

\begin{solution}{Q3. Powers of \(3\)}
Each step left to right divides by \(3\), so the indices go down by \(1\):
\[
\begin{array}{c|ccccccc}
\text{Power}&3^3&3^2&3^1&3^0&3^{-1}&3^{-2}&3^{-3}\\[2pt]
\hline
\text{Value}&27&9&3&1&\frac13&\frac19&\frac1{27}
\end{array}
\]
Thus the missing powers are \(3^1,3^{-1},3^{-2}\), and the missing values are
\[
3,\quad 1,\quad \frac19,\quad \frac1{27}.
\]
\end{solution}

\begin{solution}{Q4. Zero and negative indices}
Use \(a^0=1\) and \(a^{-n}=\dfrac1{a^n}\).
\begin{enumerate}[label=\textbf{(\alph*)}]
\item \(4^0=1\).
\item \(6^{-1}=\dfrac16\).
\item \(3^{-2}=\dfrac1{3^2}=\dfrac19\).
\item \(2^{-4}=\dfrac1{2^4}=\dfrac1{16}\).
\item \(10^{-2}=\dfrac1{10^2}=\dfrac1{100}\).
\item \(5^{-3}=\dfrac1{5^3}=\dfrac1{125}\).
\item \(1^{-7}=\dfrac1{1^7}=1\).
\item \(2^{-1}+2^{-2}=\dfrac12+\dfrac14=\dfrac24+\dfrac14=\dfrac34\).
\end{enumerate}
\end{solution}

\begin{solution}{Q5. Using the laws with negative indices}
Use \(a^m a^n=a^{m+n}\) and \(\dfrac{a^m}{a^n}=a^{m-n}\).
\begin{enumerate}[label=\textbf{(\alph*)}]
\item \(2^{-3}\times2^5=2^{-3+5}=2^2\).
\item \(7^2\div7^5=7^{2-5}=7^{-3}\).
\item \(5^4\times5^{-6}=5^{4-6}=5^{-2}\).
\item \(x^{-2}\times x^{-3}=x^{-2-3}=x^{-5}\).
\item \(3^{-2}\div3^4=3^{-2-4}=3^{-6}\).
\item \(a^3\div a^{-2}=a^{3-(-2)}=a^{3+2}=a^5\).
\item \(\dfrac{y^{-4}}{y^{-6}}=y^{-4-(-6)}=y^{-4+6}=y^2\).
\item \(m^{-1}\times m=m^{-1+1}=m^0=1\).
\item \(\dfrac{4^2}{4^{-3}\times4^8}=\dfrac{4^2}{4^{-3+8}}=\dfrac{4^2}{4^5}=4^{2-5}=4^{-3}\).
\item \(\dfrac{b^{-2}\times b^{-5}}{b^{-3}}=\dfrac{b^{-2-5}}{b^{-3}}=\dfrac{b^{-7}}{b^{-3}}=b^{-7-(-3)}=b^{-4}\).
\end{enumerate}
\end{solution}

\begin{solution}{Q6. Powers of powers}
Use \((a^m)^n=a^{mn}\).
\begin{enumerate}[label=\textbf{(\alph*)}]
\item \((5^2)^3=5^{2\times3}=5^6\).
\item \((x^4)^5=x^{4\times5}=x^{20}\).
\item \((3^4)^2=3^{4\times2}=3^8\).
\item \((y^5)^3=y^{5\times3}=y^{15}\).
\item \((10^2)^{-3}=10^{2\times(-3)}=10^{-6}\).
\item \((a^{-2})^{-4}=a^{(-2)\times(-4)}=a^8\).
\item \((7^{-1})^5=7^{-1\times5}=7^{-5}\).
\item \((6^3)^0=6^{3\times0}=6^0=1\).
\item \((p^3)^4=p^{3\times4}=p^{12}\). The missing index is \(3\).
\item \((m^{-3})^{-2}=m^{(-3)\times(-2)}=m^6\). The missing index is \(-2\).
\end{enumerate}
\end{solution}

\begin{solution}{Q7. Harder: writing numbers as powers of \(2\)}
Use \(4=2^2\), \(8=2^3\), \(16=2^4\), \(32=2^5\).
\begin{enumerate}[label=\textbf{(\alph*)}]
\item \(4^3=(2^2)^3=2^{2\times3}=2^6\).
\item \(8^4=(2^3)^4=2^{3\times4}=2^{12}\).
\item \(8^2\times4^{-1}=(2^3)^2\times(2^2)^{-1}=2^6\times2^{-2}=2^{6-2}=2^4\).
\item \(\dfrac1{32}=\dfrac1{2^5}=2^{-5}\).
\item \(16^{-2}=(2^4)^{-2}=2^{4\times(-2)}=2^{-8}\).
\item \(\dfrac{4^5}{8^3}=\dfrac{(2^2)^5}{(2^3)^3}=\dfrac{2^{10}}{2^9}=2^{10-9}=2^1\).
\end{enumerate}
\end{solution}

\begin{solution}{Q8. Simplifying expressions with indices}
Multiply or divide the coefficients using ordinary arithmetic, and use the index laws on each letter separately.
\begin{enumerate}[label=\textbf{(\alph*)}]
\item \(5x^3\times2x^4=(5\times2)x^{3+4}=10x^7\).
\item \(18p^5\div6p^2=(18\div6)p^{5-2}=3p^3\).
\item \(4y\times3y^6=(4\times3)y^{1+6}=12y^7\).
\item \(20m^8\div5m^2=(20\div5)m^{8-2}=4m^6\).
\item \(6a^2b\times2ab^3=(6\times2)a^{2+1}b^{1+3}=12a^3b^4\).
\item \(\dfrac{15x^4y^5}{3xy^2}=\dfrac{15}{3}x^{4-1}y^{5-2}=5x^3y^3\).
\item \((3x^2)^2=3^2x^{2\times2}=9x^4\).
\item \((2a^3b)^3=2^3a^{3\times3}b^{1\times3}=8a^9b^3\).
\item \(3x^{-2}\times4x^5=(3\times4)x^{-2+5}=12x^3\).
\item \(\dfrac{10y^3}{2y^5}=\dfrac{10}{2}y^{3-5}=5y^{-2}=\dfrac{5}{y^2}\).
\item \((5k^{-1})^2=5^2k^{-1\times2}=25k^{-2}=\dfrac{25}{k^2}\).
\end{enumerate}
\end{solution}

\begin{solution}{Q9. Area and volume with indices}
\begin{enumerate}[label=\textbf{(\alph*)}]
\item Rectangle:
\[
\text{Area}=\text{length}\times\text{width}=(3x^4)(2x^3)
=(3\times2)x^{4+3}=6x^7\ \text{cm}^2.
\]
\item Cube:
\[
\text{Volume}=(\text{edge})^3=(2a^3)^3
=2^3(a^3)^3=8a^9\ \text{cm}^3.
\]
\item Given area \(24x^7\) and width \(4x^2\),
\[
\text{length}=\dfrac{24x^7}{4x^2}
=\dfrac{24}{4}x^{7-2}=6x^5\ \text{cm}.
\]
\end{enumerate}
\end{solution}

\begin{solution}{Q10. Three expressions simplifying to \(16x^{12}\)}
Many answers are possible. For example:
\[
\underbrace{4x^5\times4x^7}_{\text{multiplication}}=16x^{12},\qquad
\underbrace{32x^{15}\div2x^3}_{\text{division}}=16x^{12},\qquad
\underbrace{(4x^6)^2}_{\text{bracket}}=16x^{12}.
\]
Checking:
\[
4x^5\times4x^7=(4\times4)x^{5+7}=16x^{12},
\]
\[
32x^{15}\div2x^3=(32\div2)x^{15-3}=16x^{12},
\]
\[
(4x^6)^2=4^2x^{6\times2}=16x^{12}.
\]
\end{solution}

\begin{solution}{Q11. Fractions as the base}
Use \(\left(\dfrac ab\right)^n=\dfrac{a^n}{b^n}\) and
\(\left(\dfrac ab\right)^{-n}=\left(\dfrac ba\right)^n\).
\begin{enumerate}[label=\textbf{(\alph*)}]
\item \(\left(\dfrac34\right)^2=\dfrac{3^2}{4^2}=\dfrac9{16}\).
\item \(\left(\dfrac25\right)^{-2}=\left(\dfrac52\right)^2=\dfrac{25}{4}\).
\item \(\left(\dfrac12\right)^3=\dfrac{1^3}{2^3}=\dfrac18\).
\item \(\left(\dfrac23\right)^3=\dfrac{2^3}{3^3}=\dfrac8{27}\).
\item \(\left(\dfrac14\right)^{-1}=\dfrac41=4\).
\item \(\left(\dfrac34\right)^{-1}=\dfrac43\).
\item \(\left(\dfrac13\right)^{-2}=3^2=9\).
\item \(\left(\dfrac23\right)^{-3}=\left(\dfrac32\right)^3=\dfrac{27}{8}\).
\item \(\left(\dfrac52\right)^{-2}=\left(\dfrac25\right)^2=\dfrac4{25}\).
\item \(\left(\dfrac79\right)^0=1\).
\end{enumerate}
\end{solution}

\begin{solution}{Q12. Decimals as bases}
Write each decimal as a simplified fraction before applying the index.

\textbf{(a) \(0.6^{-2}\)}
\[
0.6=\frac{6}{10}=\frac35,
\]
so
\[
0.6^{-2}=\left(\frac35\right)^{-2}
=\left(\frac53\right)^2
=\frac{5^2}{3^2}
=\frac{25}{9}.
\]
The gaps are:
\[
\frac{6}{10};\qquad \frac35;\qquad
\left(\frac35\right)^{-2};\qquad
\left(\frac53\right)^2;\qquad
\frac{25}{9}.
\]

\textbf{(b) \(0.25^{-3}\)}
\[
0.25=\frac{25}{100}=\frac14,
\]
so
\[
0.25^{-3}=\left(\frac14\right)^{-3}
=\left(\frac41\right)^3
=4^3
=64.
\]
The gaps are:
\[
\frac{25}{100};\qquad \frac14;\qquad
\left(\frac14\right)^{-3};\qquad
\left(\frac41\right)^3=4^3;\qquad
64.
\]

\textbf{(c)--(l)} Using the same method:
\begin{enumerate}[label=\textbf{(\alph*)},start=3]
\item \(0.3^2=\left(\dfrac3{10}\right)^2=\dfrac9{100}=0.09\).
\item \(0.2^3=\left(\dfrac15\right)^3=\dfrac1{125}=0.008\).
\item \(0.5^{-1}=\left(\dfrac12\right)^{-1}=2\).
\item \(0.5^{-3}=\left(\dfrac12\right)^{-3}=2^3=8\).
\item \(0.1^{-2}=\left(\dfrac1{10}\right)^{-2}=10^2=100\).
\item \(0.2^{-2}=\left(\dfrac15\right)^{-2}=5^2=25\).
\item \(1.5^2=\left(\dfrac32\right)^2=\dfrac94=2.25\).
\item \(0.75^{-2}=\left(\dfrac34\right)^{-2}=\left(\dfrac43\right)^2=\dfrac{16}{9}\).
\item \(2.5^{-2}=\left(\dfrac52\right)^{-2}=\left(\dfrac25\right)^2=\dfrac4{25}=0.16\).
\item \(0.8^{-1}=\left(\dfrac45\right)^{-1}=\dfrac54=1.25\).
\end{enumerate}
\end{solution}

\begin{solution}{Q13. Using the laws, then working out the value}
\begin{enumerate}[label=\textbf{(\alph*)}]
\item \(0.5^3\times0.5^{-5}=0.5^{3-5}=0.5^{-2}
=\left(\dfrac12\right)^{-2}=2^2=4\).
\item \(\left(\dfrac23\right)^5\div\left(\dfrac23\right)^7
=\left(\dfrac23\right)^{5-7}
=\left(\dfrac23\right)^{-2}
=\left(\dfrac32\right)^2=\dfrac94\).
\item \((0.2^{-1})^2=0.2^{-2}
=\left(\dfrac15\right)^{-2}=5^2=25\).
\item \(\dfrac{0.1^4}{0.1^7}=0.1^{4-7}=0.1^{-3}
=\left(\dfrac1{10}\right)^{-3}=10^3=1000\).
\item \(0.4^{-1}\times0.4^{-1}=0.4^{-2}
=\left(\dfrac25\right)^{-2}
=\left(\dfrac52\right)^2=\dfrac{25}{4}\).
\item \((0.25^2)^{-1}=0.25^{-2}
=\left(\dfrac14\right)^{-2}=4^2=16\).
\end{enumerate}
\end{solution}

\begin{solution}{Q14. Tiles and cubes}
\begin{enumerate}[label=\textbf{(\alph*)}]
\item
\[
0.25^{-2}=\left(\frac14\right)^{-2}=4^2=16.
\]
There are \(4\) tiles along each side of the \(1\) m floor, so the floor contains
\[
4\times4=16
\]
tiles. This agrees with \(0.25^{-2}=16\).

\item A small cube has edge \(0.2\) m, so its volume is \(0.2^3\) m\(^3\). The number of small cubes in the \(1\) m\(^3\) crate is
\[
\frac{1}{0.2^3}=0.2^{-3}
=\left(\frac15\right)^{-3}
=5^3
=125.
\]
So \(125\) small cubes fit in the crate.
\end{enumerate}
\end{solution}

\begin{solution}{Q15. Reasoning questions}
\begin{enumerate}[label=\textbf{(\alph*)}]
\item Sam is not correct. A negative index means take the reciprocal of the base:
\[
0.5^{-2}=\left(\frac12\right)^{-2}=2^2=4,
\]
which is positive. The negative index does not make the value negative.

\item Work out each number:
\[
0.5^{-3}=2^3=8,\qquad
0.2^{-2}=5^2=25,\qquad
\left(\frac32\right)^{-2}=\left(\frac23\right)^2=\frac49,
\]
\[
0.1^{-1}=10,\qquad
4^{-1}=\frac14.
\]
From smallest to largest:
\[
4^{-1}=\frac14
\;<\;
\left(\frac32\right)^{-2}=\frac49
\;<\;
0.5^{-3}=8
\;<\;
0.1^{-1}=10
\;<\;
0.2^{-2}=25.
\]

\item Many answers are possible. For example:
\[
16=4^2=16^1=0.25^{-2}
=\left(\frac14\right)^{-2}
=\left(\frac1{16}\right)^{-1}
=0.0625^{-1}
=(-4)^2
=(-2)^4.
\]
Also \(16=2^4\) and \(16=0.5^{-4}\), as given in the question.
\end{enumerate}
\end{solution}

\end{document}